\documentclass[a4paper]{article}

\usepackage{authblk}

\usepackage{graphicx}
\graphicspath{{figures/}}

\usepackage{hyperref}

\hypersetup{
  pdfauthor = {Nikolaos Pothitos},
  pdftitle = {Naxos Solver}
}

\begin{document}

\title{\begin{picture}(0,0)%
\includegraphics{logo}%
%\special{psfile=logo/naxos}%
\end{picture}%
\setlength{\unitlength}{4144sp}%
%
\begingroup\makeatletter\ifx\SetFigFontNaxos\undefined
% extract first six characters in \fmtname
\def\x#1#2#3#4#5#6#7\relax{\def\x{#1#2#3#4#5#6}}%
\expandafter\x\fmtname xxxxxx\relax \def\y{splain}%
\ifx\x\y   % LaTeX or SliTeX?
\gdef\SetFigFontNaxos#1#2#3{%
  \ifnum #1<17\tiny\else \ifnum #1<20\small\else
  \ifnum #1<24\normalsize\else \ifnum #1<29\large\else
  \ifnum #1<34\Large\else \ifnum #1<41\LARGE\else
     \huge\fi\fi\fi\fi\fi\fi
  \csname #3\endcsname}%
\else
\gdef\SetFigFontNaxos#1#2#3{\begingroup
  \count@#1\relax \ifnum 25<\count@\count@25\fi
  \def\x{\endgroup\@setsize\SetFigFontNaxos{#2pt}}%
  \expandafter\x
    \csname \romannumeral\the\count@ pt\expandafter\endcsname
    \csname @\romannumeral\the\count@ pt\endcsname
  \csname #3\endcsname}%
\fi
\fi\endgroup
\begin{picture}(2166,2991)(-32,-2344)
\put(901,-1411){\makebox(0,0)[lb]{\smash{\SetFigFontNaxos{20}{24.0}{rm}\textsc{\LARGE\ Naxos\ \ Solver}}}}
\end{picture}
}

\author{Nikolaos Pothitos}

\affil{\normalsize
       Department of Informatics and Telecommunications \\
       National and Kapodistrian University of Athens \\
       Panepistimiopolis, 157\,84 Athens, Greece \\
       \texttt{pothitos@di.uoa.gr}}

\date{}

\maketitle

\begin{itemize}
  \item \textsc{Naxos Solver} is a C++ Constraint
        Programming library.
  \item It is freely available in the open source community
        together with its full documentation~\cite{Naxos}.
  \item It enforces a bounds-consistency variant in order to
        propagate changes that happen in Constraint
        Satisfaction Problem variables~\cite{Bessiere2006}.
        This happens in every single node of the search
        tree.
  \item It supports the definition of custom search methods
        via a declarative API~\cite{Pothitos2017}.
  \item Declarativity is a \textsc{Naxos} key feature and it
        has do with describing problems and search methods
        independently. \textsc{Naxos} serves the Constraint
        Programming principle: \emph{The user states the
        problem, the computer solves it}~\cite{Freuder2014}.
  \item \textsc{Naxos} is being used in various
        applications, such as the construction of the
        timetable for the Department of Informatics and
        Telecommunications of the University of Athens every
        semester~\cite{Pothitos2012-Scheduling}.
\end{itemize}

\section*{Acknowledgments}

I would like to thank Prof.~Panagiotis Stamatopoulos for his
inspiration to build the solver and his continuous
supervision and scientific support. I also thank Foivos
Theocharis for the search methods modules implementation.

\bibliographystyle{abbrv}
\bibliography{bibliography}

\end{document}
